\documentclass[11pt]{ctexart}
\usepackage[a4paper,top=2.2cm,bottom=2.2cm,left=2.2cm,right=2.2cm]{geometry}
\usepackage{booktabs}
\usepackage{longtable}
\usepackage{array}
\usepackage{enumitem}
\usepackage{xcolor}
\usepackage{amssymb}
\usepackage{caption}
\usepackage[hidelinks]{hyperref}
\setlist{nosep, leftmargin=1.6em}
\definecolor{grayline}{gray}{0.65}
\emergencystretch=1.5em
\setlength{\tabcolsep}{4pt}
\renewcommand{\arraystretch}{1.25}
\counterwithin{table}{section}
\counterwithin{figure}{section}
\renewcommand{\thetable}{\thesection-\arabic{table}}
\renewcommand{\thefigure}{\thesection-\arabic{figure}}

\ctexset{
  section = {name = {,、}, number = \chinese{section}, format = \heiti\zihao{3},
             beforeskip = 1.2ex plus .2ex minus .2ex, afterskip = .6ex plus .1ex minus .1ex},
  subsection = {name = {（,）}, number = \chinese{subsection}, format = \heiti\zihao{4},
                beforeskip = .9ex plus .1ex minus .1ex, afterskip = .4ex plus .1ex minus .1ex},
}
\captionsetup{labelsep=quad, font=small, skip=6pt}

\hypersetup{
  pdftitle={SmartFarmAgent 智慧农业智能体系统——参赛作品说明书},
  pdfsubject={参赛作品说明书（v3.3 开源资源与赛事链接版）},
  pdfauthor={}
}

\begin{document}

\begin{center}
  {\LARGE\bfseries 参赛作品说明书}\\[8pt]
  {\large\bfseries SmartFarmAgent——基于大模型智能体的设施农业「感知—决策—执行」闭环系统}\\[4pt]
  {\normalsize 完成时间：2026 年 8 月}
\end{center}

\vspace{1.2em}

\section{作品简介}
\subsection{背景与痛点}
设施农业（大棚、连栋温室）是保障果蔬稳定供给的重要生产方式，但传统管理高度依赖人工巡检与经验调控：浇水靠手感、病害靠肉眼、生长判断靠主观，存在缺水涝害、病害发现滞后、水肥浪费、环境调控不及时等问题。现有市售「智慧农业」产品大多采用阈值规则控制（如土壤湿度低于设定值即浇水），缺少多源数据综合判断能力；大模型应用多停留在「知识问答」层，只回答问题、不驱动设备，难以形成真正的智能闭环。企业级智慧农业平台功能完整，但重软件集成、成本高、部署复杂，不适合小规模设施农业与教学科研场景快速落地。

\subsection{行业与赛事趋势}
农业农村部信息中心等单位于 2026 年 7 月在雄安举办智慧农业创新大赛，设智能农机作业控制、无人机巡田、激光除草、设施小番茄采摘、家禽巡检、鱼群智能检测六大赛道，采用实景场地、真机实操、现场竞技，重点检验装备在真实生产场景中的作业能力；同期国际大学生智能农业装备创新大赛（天鹅杯）、全国大学生物联网设计竞赛（华为杯）、挑战杯「人工智能+」专项赛、中国国际大学生创新大赛等获奖作品普遍具备「软硬一体化、感知—决策—执行闭环、量化效果、场景验证」的共性。低成本、可复现、闭环落地的软硬一体系统，是当前竞赛与行业共同认可的立项方向。

\subsection{作品定位}
本作品以 ESP32 单芯片为感知执行端，Node.js + Vue 为云平台，大模型智能体为决策大脑，构建设施农业「感知—决策—执行」全闭环系统：传感器与摄像头持续采集环境与作物状态，云端智能体综合传感器、视觉与历史数据给出设备指令和农技方案，指令经白名单校验后下发执行并回传状态，形成数据驱动、可追溯的完整业务闭环。

\subsection{核心理念}
数据驱动、检索前置、本地优先、闭环落地。本地阈值规则永远优先于云端指令（断网自持、误操作兜底），云端大模型指令走白名单；软硬一体、低成本单芯片，整套系统可快速复现，适配教学科研、小规模经营与竞赛展示。

\section{设计原理}
\subsection{系统分层架构}
系统按六层组织，各层职责与关键技术如表 \ref{tab:layers} 所示。
\begin{longtable}{p{1.2cm}p{2.2cm}p{8.8cm}p{3.0cm}}
  \caption{系统分层架构}\label{tab:layers}\\
  \toprule
  层级 & 名称 & 组成与职责 & 关键技术 \\
  \midrule
  L6 & 展示层 & Vue 3 Web 大屏 / H5：实时曲线、设备控制、农技问答、视觉告警 & Vue、ECharts、WebSocket \\
  L5 & 智能决策层 & 大模型智能体（LLM+RAG）→ JSON 设备指令；视觉服务（YOLO 病害/生长分析） & LLM、RAG、视觉 \\
  L4 & 云端平台层 & Node.js 后端：MQTT 接入 → 时序存储 → REST/WebSocket API；智能体编排 & Node.js、MQTT、时序存储 \\
  L3 & 边缘层 & ESP32-CAM 定时拍照上传；断网规则兜底（本地优先） & 边缘计算、定时任务 \\
  L2 & 通信层 & ESP32 内置 Wi-Fi：MQTT（JSON 上行/下行）+ HTTP 图片上传 & Wi-Fi、MQTT \\
  L1 & 感知执行层 & ESP32 单芯片：传感器采集 + 继电器/泵/灯/喷药 + OLED + microSD 日志 & GPIO、ADC、I2C、PWM \\
  \bottomrule
\end{longtable}
安全原则：本地阈值规则永远优先于云端指令（断网自持、误操作兜底），云端智能体指令走白名单。

\subsection{核心业务模块}
\begin{enumerate}
  \item \textbf{感知采集与本地日志}：BH1750 光照、DHT11 温湿度、土壤湿度等传感器由 ESP32 定时采集，打包为 JSON 上行；同时写入 microSD 本地日志，断网时数据不丢失，联网后补传。
  \item \textbf{环境自动调控}：根据温湿度与光照综合判断，驱动风扇、加热片与补光灯，PWM 调光调速（PID 可选），实现设施环境自动调节。
  \item \textbf{病虫害检测与给药}：ESP32-CAM 定时拍照上传，YOLO 模型检测叶片病害，置信度达标后自动触发喷药泵，实现「识别—决策—执行」联动。
  \item \textbf{生长状况分析}：定距定时拍照，统计叶片数、株高、颜色等表型信息，由大模型生成阶段性生长周报，形成可量化的生长记录。
  \item \textbf{农技知识问答}：种植、病虫害、用药知识库经 RAG 检索增强生成结构化方案，回答附来源标注，抑制模型幻觉，可溯源。
  \item \textbf{智能体决策}：Node.js 编排大模型智能体，综合传感器、视觉与历史数据输出 JSON 设备指令与解释；指令经白名单校验后下发 ESP32 执行，误操作有兜底。
  \item \textbf{数据中台与可视化}：MQTT 上行数据经 Node.js 时序存储，Vue + ECharts 大屏实时展示；执行状态回传形成下一次决策依据。
\end{enumerate}

\subsection{大模型与智能体协同机制}
系统按「检索前置、模型后置、数据驱动」运行：RAG 从真实农业知识库构建上下文，大模型据此生成回答与方案；智能体把方案转换为 JSON 设备指令，经白名单校验后下发执行，执行结果回传闭环。纯软件问答系统只回答、不驱动设备，纯阈值规则控制缺少多源综合判断；本作品把两者接入同一条链路，既保留自然语言交互，也能把决策落到设备动作。

\subsection{完整数据链路}
数据从感知到执行、再回到决策的完整链路如表 \ref{tab:datapath} 所示。
\begin{longtable}{p{1.5cm}p{11.0cm}p{2.6cm}}
  \caption{完整数据链路}\label{tab:datapath}\\
  \toprule
  环节 & 内容 & 输出 \\
  \midrule
  感知 & BH1750/DHT11/土壤 → ESP32 采集打包（JSON）；ESP32-CAM 定时拍照 & 数据 / 图片 \\
  通信 & ESP32 Wi-Fi：MQTT 上行情报 + HTTP 图片上传；断网时 microSD 本地日志 & 上行数据 / 图片 \\
  云端 & Node.js 后端（mqtt.js）订阅 → 时序存储 → REST/WebSocket API & 时序数据 \\
  AI & 视觉服务（YOLO 病害检测/生长分析）→ 结果入事件队列；智能体（LLM+RAG）综合传感器+视觉+历史 & 决策 JSON / 问答方案 \\
  执行 & MQTT 下行指令 → ESP32 解析（白名单）→ 继电器/泵/灯/喷药 & 设备动作 \\
  闭环 & 执行状态回传 → 数据入库 → Vue 大屏展示 + 智能体下次决策依据 & 反馈数据 \\
  \bottomrule
\end{longtable}

\section{技术方案}
\subsection{技术栈}
各技术层选型与作用如表 \ref{tab:stack} 所示。
\begin{longtable}{p{2.4cm}p{4.8cm}p{7.6cm}}
  \caption{技术栈}\label{tab:stack}\\
  \toprule
  技术层 & 采用技术 & 作用说明 \\
  \midrule
  前端界面 & Vue 3、ECharts、WebSocket & 大屏看板、设备控制、农技问答、视觉告警等交互界面 \\
  后端服务 & Node.js、Express、MQTT.js & 业务逻辑、接口路由、MQTT 接入、智能体编排与权限校验 \\
  通信协议 & MQTT（JSON）/ HTTP / WebSocket & 设备上行数据、下行指令与图片上传、前端实时推送 \\
  数据存储 & MySQL / SQLite 时序存储 & 传感器数据、设备状态、问答会话与配置等结构化数据 \\
  视觉服务 & Python、OpenCV、YOLOv8 & 病害检测、生长表型分析（独立 AI 服务，由 Node.js 编排调用） \\
  大模型与 RAG & 兼容 OpenAI API 的模型服务（如 DeepSeek）、RAGFlow/Dify 或自建向量检索 & 农技问答、智能体决策与知识溯源 \\
  边缘端 & ESP32（Arduino/ESP-IDF）、ESP32-CAM & 数据采集、拍照、本地规则控制、microSD 日志 \\
  \bottomrule
\end{longtable}

\subsection{硬件方案}
主要器件清单与用途如表 \ref{tab:hardware} 所示。
\begin{longtable}{p{6.4cm}p{8.8cm}}
  \caption{硬件方案}\label{tab:hardware}\\
  \toprule
  器件 & 用途 \\
  \midrule
  ESP32-CAM 开发板（OV2640 + microSD 卡槽） & 主控：数据采集、拍照、联网、控制 \\
  BH1750 光照传感器 & 光照采集（I2C） \\
  DHT11 温湿度传感器 & 空气温湿度采集 \\
  土壤湿度传感器 & 土壤水分采集（ADC） \\
  OLED 显示屏 & 本地数据显示（I2C） \\
  继电器模块 & 设备开关控制 \\
  水泵 / 蠕动泵 & 浇水、施肥/给药 \\
  风扇 / 加热片 / 补光灯 & 通风、低温补偿、光照补充（PWM 调光） \\
  5V 稳压供电模块 & ESP32-CAM 稳定供电（电流敏感） \\
  microSD 卡 & 断网本地日志存储 \\
  \bottomrule
\end{longtable}

\section{创新点}
\subsection{低成本单芯片软硬一体闭环}
采用 ESP32 单芯片完成采集、拍照、控制、联网与日志存储，链路最短、硬件总成本低，整套「感知—决策—执行」闭环可在千元级硬件上完整复现，突破传统方案软硬分离、成本高的局限。

\subsection{大模型决策闭环到设备动作}
大模型智能体不仅回答问题，还直接输出 JSON 设备指令；指令经白名单校验、本地规则复核后下发执行，兼顾智能化与安全性。这是与纯软件问答系统、纯阈值控制系统的核心区别。

\subsection{边缘—云协同与断网自持}
断网时本地阈值规则独立运行并写入 microSD，联网后数据补传、云端智能体接管；本地规则永远优先，关键动作可设人工确认模式，误操作有兜底。

\subsection{农业 RAG 溯源问答}
种植与植保知识库经检索增强生成回答，附来源标注，抑制大模型幻觉，使农技建议可查证、可追溯，适合严肃农业场景。

\subsection{全周期生长数字档案}
从定植到收获持续记录传感器与视觉数据，生成可追溯、可对比的生长档案，支撑产量预测与栽培方案优化，形成数据资产。

\subsection{创新效果对比}
创新点、实现方式与量化目标的对比如表 \ref{tab:innovation} 所示。
\begin{longtable}{p{2.3cm}p{4.4cm}p{4.2cm}p{3.6cm}}
  \caption{创新效果对比}\label{tab:innovation}\\
  \toprule
  创新点 & 技术方案 & 与现有方法对比 & 量化目标 \\
  \midrule
  单芯片软硬一体 & ESP32 单芯片集成采集/拍照/控制/联网 & 传统方案软硬分离、成本高、部署繁琐 & 硬件总成本千元级，5 分钟现场部署 \\
  智能体闭环决策 & LLM+RAG → JSON 指令 → 白名单 → 执行回传 & 纯问答无执行、纯规则无智能 & 指令闭环响应 < 2 秒，误动作拦截率 100\% \\
  边缘-云协同 & 本地规则优先 + microSD 断网补传 & 云端依赖型方案断网即失控 & 断网自持 $\geq$ 24 小时，日志零丢失 \\
  RAG 溯源问答 & 知识库检索增强 + 来源标注 & 通用大模型易幻觉、不可溯源 & 问答溯源覆盖率 100\%，关键指标可核验 \\
  生长数字档案 & 传感器+视觉持续记录与对比 & 人工记录零散、不可追溯 & 全周期档案连续可查，周报自动生成 \\
  \bottomrule
\end{longtable}

\section{实用价值与应用场景}
\subsection{降低管理门槛、提升响应效率}
管理者以自然语言提问或下达指令即可完成查询与调控，无需学习数据库与编程；环境异常自动告警，病害早发现早处置，减少水肥浪费与产量损失。

\subsection{数据真实可靠、可追溯}
所有回答与方案均基于本地知识库与实测数据生成，附来源标注；设备动作全记录，执行状态可回查，满足严肃生产场景对可靠性的要求。

\subsection{低成本可复制}
千元级硬件 + 开源软件栈，单套系统可快速复现，适合家庭农场、教学实验、科研平台与竞赛展示等低成本场景。

\subsection{模块化可推广}
感知与执行模块可替换：更换传感器与执行器即可迁移到水产养殖、禽舍巡检、果园管理、大田监测等场景；平台与智能体底座复用，推广性强。

\subsection{竞赛与展示适配}
系统具备真机实操演示能力：现场可展示传感器实时曲线、自然语言控制设备、病害识别联动喷药、断网自持等完整闭环，配套说明书、演示视频、易拉宝等材料，适配多类赛事评审要求。

\section{总结}
本作品以「感知—决策—执行」闭环为核心，采用 ESP32 单芯片、Node.js/Vue 平台与云端大模型智能体，将数据采集、视觉识别、智能决策与设备执行打通为一个可运行、可复现、可推广的软硬一体系统。通过本地规则优先与指令白名单双重安全机制、RAG 溯源问答与全周期数字档案，解决传统设施农业管理依赖人工、智能决策不闭环、数据不可追溯等痛点。系统以低成本单芯片实现完整闭环，紧贴当前全国性竞赛「实景实操、闭环落地、量化验证」的评审导向，具备良好的竞赛适配性与成果转化前景。

\appendix
\ctexset{section={name={附录~,：}, number=\Alph{section}, format=\heiti\zihao{3}}}
\section{赛事适配说明}
全国 9 类相关赛事已完成官网核验并按含金量排序，明细见独立文档《全国相关赛事调研报告》（Competition\_Research，同目录附 MD/PDF/DOCX）。核心结论：
\begin{itemize}
  \item \textbf{技术栈主攻}：\href{https://iot.sjtu.edu.cn}{全国大学生物联网设计竞赛（华为杯）}、\href{https://www.socchina.net}{全国大学生嵌入式芯片与系统设计竞赛}、\href{https://www.tiaozhanbei.net}{挑战杯「人工智能+」专项赛}。
  \item \textbf{行业主攻}：\href{https://www.moa.gov.cn/xw/bmdt/202605/t20260520_6484332.htm}{智慧农业创新大赛}（农业农村部信息中心）、\href{http://uiaec.ujs.edu.cn}{国际大学生智能农业装备创新大赛（天鹅杯）}。
  \item \textbf{获奖共性}：实景实操、软硬一体、大模型/智能体落到执行、量化验证与场景验证（示例：\href{https://www.moa.gov.cn/xw/bmdt/202605/t20260520_6484332.htm}{2026 智慧农业创新大赛}设施小番茄赛道满分 200 分；\href{https://www.tiaozhanbei.net}{挑战杯「人工智能+」}一等奖「花语智农」）。
\end{itemize}
申报策略：以「感知—决策—执行闭环能力底座」统一叙事，按赛事切换侧重点（物联网/嵌入式类重工程实现，挑战杯/创新大赛重 AI 应用与社会价值，农业装备类重装备化与作业指标）；材料按「作品说明书 + 技术文档 + 演示视频 + 易拉宝 + 承诺书」标准包准备。
\section{里程碑计划（约 6 个月）}
六个阶段的安排与验收点如表 \ref{tab:milestone} 所示。
\begin{longtable}{p{2.2cm}p{1.7cm}p{6.5cm}p{4.0cm}}
  \caption{里程碑计划}\label{tab:milestone}\\
  \toprule
  阶段 & 时间 & 内容 & 验收点 \\
  \midrule
  M0 开题 & 第 1—2 周 & 链路定稿（ESP32 + Node.js + Vue）、软件环境搭建 & 开题通过 \\
  M1 感知-控制闭环 & 第 3—6 周 & ESP32 传感器采集 → OLED → 继电器 → MQTT（本地 EMQX）→ Node.js → Vue 控制页 & 局域网内数据正确、远程手动可控 \\
  M2 上云+拍照+日志 & 第 7—10 周 & ESP32-CAM 定时拍照上传；MQTT 上云；Node.js 时序存储 + Vue 大屏；microSD 断网补传 & 实时曲线、图片可查、断网日志完整可补传 \\
  M3 视觉检测 & 第 11—18 周 & 公开集 + 自采数据训练 YOLO（先 5 类常见病）→ Node.js 调视觉服务 → 触发喷药 & 检测准确率 $\geq$ 85\%，喷药动作联动 \\
  M4 智能体决策 & 第 19—23 周 & Node.js 编排智能体 + RAG 知识库（设备决策 + 农技问答）→ JSON 指令闭环 → Vue 对话 & 自然语言指令自动执行，问答可溯源，误操作有兜底 \\
  M5 作品化 & 第 24 周 & 外壳/说明书/演示视频/易拉宝/答辩材料 & 达到参赛交付标准 \\
  \bottomrule
\end{longtable}

\section{风险与应对}
主要风险与应对措施如表 \ref{tab:risk} 所示。
\begin{longtable}{p{5.2cm}p{10.0cm}}
  \caption{风险与应对}\label{tab:risk}\\
  \toprule
  风险 & 应对 \\
  \midrule
  视觉识别数据不足/效果差 & 公开数据集预训练 + 自采微调；范围先缩到 3—5 类常见病害 \\
  ESP32-CAM 算力不足 & 初期视觉推理放云端，边缘只拍照上传；后期再试轻量模型 \\
  ESP32-CAM 引脚少/供电敏感 & 稳压供电；继电器组用 I/O 扩展（如 PCF8574/74HC595）或精简执行器数量 \\
  智能体误操作 & 本地规则优先、指令白名单、关键动作人工确认模式 \\
  Wi-Fi 不稳定 & 本地阈值模式 + microSD 日志补传 \\
  时间不足 & 裁剪顺序：合并施肥/给药、生长分析降级为拍照+周报、Vue 先做单页大屏 \\
  \bottomrule
\end{longtable}

\section{参考开源资源}
系统核心链路所依赖的主要开源资源如表 \ref{tab:oss} 所示，全部为开发与部署环节直接使用的成熟项目。
\begin{longtable}{p{3.6cm}p{5.8cm}p{5.8cm}}
  \caption{参考开源资源}\label{tab:oss}\\
  \toprule
  开源资源 & 用途 & 地址 \\
  \midrule
  arduino-esp32 & 感知执行端固件开发（ESP32 官方 Arduino 核心） & \url{https://github.com/espressif/arduino-esp32} \\
  Express & 云端平台后端 Web 框架 & \url{https://github.com/expressjs/express} \\
  MQTT.js & 设备上下行 MQTT 客户端库 & \url{https://github.com/mqttjs/MQTT.js} \\
  Vue 3 & 网页大屏与交互界面框架 & \url{https://github.com/vuejs/core} \\
  Apache ECharts & 实时曲线与数据可视化 & \url{https://github.com/apache/echarts} \\
  Ultralytics YOLOv8 & 病害检测与生长分析模型 & \url{https://github.com/ultralytics/ultralytics} \\
  OpenCV & 图像预处理与分析 & \url{https://github.com/opencv/opencv} \\
  EMQX & 本地 MQTT 消息服务器 & \url{https://github.com/emqx/emqx} \\
  RAGFlow & 农技知识库检索增强（RAG）引擎 & \url{https://github.com/infiniflow/ragflow} \\
  PlantVillage-Dataset & 植物病害图像数据集（模型训练） & \url{https://github.com/spMohanty/PlantVillage-Dataset} \\
  \bottomrule
\end{longtable}

\end{document}
